\subsection{问题分析：日内合同修正的经济价值}
问题二的结果表明，在仅有日前固定合同、不允许日内调整的约束下，MPC的跨时段优化增益有限，其收益大部分被负荷与光伏预测不确定性抵消。在此基础上，问题三引入日内购电合同修正机制。微网在日前市场依据光伏功率预测签订购电合同，得到计划购电量$p_t$；受光伏发电随机性与预测误差影响，日内光伏实际出力与日前预测存在显著偏差。若不允许对日前购电合同做任何修改，功率偏差只能依赖储能充放电平抑，储能触及SOC上下限后剩余偏差只能通过高价应急购电、闲置合同惩罚与弃光三类兜底手段承受，显著推高运行成本。

问题三要求在日内运行阶段引入合同修正机制：利用更新后的光伏预报，在合理约束内对尚未执行时段的购电合同适度修正，得到实际执行的最终购电合同$q_t$，以此削弱预测误差的不利影响。题目在每日0:00、6:00、12:00、18:00四个时点发布光伏预报，每个时段的最终购电量只相对0:00原计划结算一次，中间更新仅改变尚未执行时段的购电量。

为定量评估日内调账机制的经济价值，设置两组对照方案：主策略开启全部三个日内预报时点的合同调整；基线方案$\mathcal{S}=\varnothing$不使用任何日内新预报，仍保留相同的滚动调度方法与结算口径。两组方案采用完全一致的负荷、光伏时序、储能参数、价格体系与预测输入，仿真区间为2025-02-01至2025-12-31，共计334天。

\subsection{决策时序与信息边界}
一天内共有四个预报版本，新预测只能在规定发布时间获得。0:00版本用于形成初始计划购电量$p_t$；6:00、12:00和18:00时，新版本只影响尚未执行的时段，已执行时段不再调整。每次预报更新先读取上一时段结束后的实际储能SOC，再调整剩余时段的购电量，最后进入下一时段的滚动调度。

\subsection{官方预报校正与分层风险}
对目标日$d$、发布时刻$r$和提前期$h$，使用历史同版本误差对官方预报校正：
\begin{equation}
\hat{V}_{d,r,h}^{\mathrm{cal}}=\max\left(0,\ \hat{V}_{d,r,h}^{\mathrm{raw}}+\mathrm{mean}_{j<d}\,e_{j,r,h}^{V}\right),
\label{eq:q3_cal}
\end{equation}
式中$e_{j,r,h}^{V}$为历史第$j$日同发布时刻、同提前期的光伏预报残差。校正后插值到10分钟网格。不同提前期的预报误差存在差异，预测时域越长通常不确定性越大。若对每个提前期单独估计分位，会因样本不足引入较大统计噪声。本文按提前期将误差分为六层：0-1、1-2、2-3、3-4、4-5、5-6小时，提前期超过6小时的远期时段统一归入第六层；0:00版本覆盖的6-24小时部分及后续版本中仍超过6小时的远期时段均按第六层处理。令$s_h$为提前期$h$对应残差的经验均方根尺度，标准化残差$z_{j,r,h}=e_{j,r,h}^{V}/s_h$，则风险裕量写为
\begin{equation}
R_{d,r,h}=s_h\,Q_\alpha\left\{z_{j,r,h}:\ d-W\le j<d\right\},\quad W=35,\ \alpha=0.77,
\label{eq:q3_risk}
\end{equation}
式中$Q_\alpha$为经验分位函数。该方法保留"越远期误差通常越大"的结构，同时避免小样本下逐层估计的不稳定。校正后1小时、6小时的误差MAE分别约为19.38 kW和92.71 kW，其增长趋势为分层风险设置提供了直接数据依据。

\begin{figure}[htbp]
\centering
\includegraphics[width=0.85\textwidth]{figures/q3-forecast-quality.png}
\caption{问题三不同预报提前期的校正质量。}
\label{fig:q3_forecast_quality}
\end{figure}

\subsection{优化模型构建：MPC滚动求解与合同调整规则}
\subsubsection{目标函数与结算规则}
在0:00形成初始计划购电量$p_t$，6:00、12:00、18:00保持已执行时段不变，以当时实际SOC重新计算剩余时段购电量。令
\begin{equation}
q_t=p_t+\delta_t^+-\delta_t^-,\quad \delta_t^+,\delta_t^-\ge0,
\label{eq:q3_qdef}
\end{equation}
式中$\delta_t^+,\delta_t^-$分别为上调、下调量。剩余时域优化目标为
\begin{equation}
\min \sum_{t\in\tau_r}\left(c_tp_t+1.5c_t\delta_t^+-0.5c_t\delta_t^-+5c_tu_t\right)-\lambda_r(S_{145}-S_{\min}),
\label{eq:q3_obj}
\end{equation}
式中$c_t$为$t$时段购电单价，$\tau_r$为第$r$次更新时刻尚未执行的时段集合，$\lambda_r$为终端储能价值罚项权重。$p_t$是0:00签订的原始合同量，$q_t$是最后一次调整后实际生效的合同量。$u_t$为规划层对未来缺口的应急购电代理变量，用于避免合同调整层无成本保留缺口；真实结算仍采用实际执行后的应急购电量。最终调整购电量始终相对0:00计划购电量结算一次：

与问题二相同，上述合同调整层中的$x_t,y_t$仅用于估计终端价值$\lambda_r(S_{145}-S_{\min})$的规划代理值，不是物理执行值。物理执行仍由第\ref{subsec:q2-model}节定义的同一滚动 LP 执行器给出：该执行器在合同调整层每次更新$q_t$后，以最新合同重新执行，保证规划层与执行层在合同口径上一致；执行层保留$\varepsilon(x_\tau+y_\tau)$项以抑制同时充放电，适用命题 3 的条件性保证。

\begin{equation}
C_{3,t}=c_t\min(p_t,q_t)+1.5c_t(q_t-p_t)^++0.5c_t(p_t-q_t)^++5c_tu_t,
\label{eq:q3_settle}
\end{equation}
式中$u_t$为真实负荷与光伏到达并完成物理执行后计算的应急购电量，不参与此前决策时刻的未来真值输入。该公式具有清楚的边际含义：下调部分按0.5$c_t$计费，上调部分按1.5$c_t$计费，电量缺口执行5倍应急购电价格。每个时段的最终购电量只相对0:00原计划结算一次，中间多次滚动优化的目标值不直接相加，因此不会重复计费。

\subsubsection{约束条件}
\begin{enumerate}
\item 功率平衡约束
\begin{equation}
q_t-w_t+v_t+y_t+u_t=L_t+x_t,
\label{eq:q3_power_balance}
\end{equation}
式中各变量均为单个10分钟时段电量（kWh）；$v_t$为光伏消纳量，$w_t$为闲置合同电量，$x_t,y_t$为充、放电量。
\item 储能系统运行约束
\begin{equation}
\begin{cases}
S_{t+1}=S_t+0.9x_t-\dfrac{y_t}{0.9}\\
S_{\min}\le S_t\le S_{\max}\\
0\le x_t,y_t\le B_{\max}
\end{cases}
\label{eq:q3_bess}
\end{equation}
式中$S_t$为储能$t$时刻储能量，$x_t,y_t$分别为该10分钟时段的充、放电电量；其上限由$P_{\max}\Delta t$给出，$S_{\min},S_{\max}$为SOC上下限。
\item 合同调整仅作用于尚未执行时段，且$q_t\ge0$；已执行时段满足$q_t=p_t$，未执行时段满足$q_t=p_t+\delta_t^+-\delta_t^-$，并且$0\le w_t\le q_t$。
\item 非负约束
\begin{equation}
u_t\ge0,\quad P_{\mathrm{curt},t}\ge0.
\label{eq:q3_nonneg}
\end{equation}
\end{enumerate}

\subsubsection{滚动求解配置}
在相邻预报发布时点之间，滚动LP同时考虑未来价格、修正后的净负荷和SOC边界，每次只执行当前时段动作，再用最新状态重新优化，构成MPC执行过程\cite{parisio2016stochasticmpc}。对$H=1,3,6,12$的全年试验中，$H=1$费用最低；$H=6$仅高4273.26元（$0.03\%$），但能够联合安排未来1小时的储能动作。综合样本费用、多时段协调能力和计算规模，主方案取$H=6$。

\subsection{仿真结果与分析}
\subsubsection{主策略与基线方案对比}
将两组方案按日聚合统计，得到334天仿真周期内的年度关键指标，如表\ref{tab:q3_compare}所示。
\begin{table}[htbp]
\centering
\caption{问题三主策略与基线方案年度指标对比}
\label{tab:q3_compare}
\scriptsize
\begin{tabular}{lccc}
\toprule
指标 & 主策略(启用全部日内更新) & 基线方案$\mathcal{S}=\varnothing$ & 变化量 \\ 
\midrule
年度总费用/元 & 14664345.40 & 15230376.12 & $-566030.72$ \\
应急购电量/kWh & 210761.48 & 334601.06 & $-123839.58$ \\
\bottomrule
\end{tabular}
\end{table}
主策略年度费用较基线下降\textbf{566030.72元}，降幅约$3.72\%$，应急购电量同时下降$37\%$。此外，0:00计划购电量为\textbf{21702724.37 kWh}，最终调整购电量为\textbf{21702916.77 kWh}；累计上调量\textbf{513849.94 kWh}、下调量\textbf{513657.53 kWh}，两者基本相抵但发生时段不同，说明合同调整将风险在各时段之间重新分配，整日合同没有单向变化；弃购量为\textbf{1228166.90 kWh}。

\subsubsection{预报信息消融实验}
为隔离日内预报信息本身的价值，在保持同一滚动调度方法、同一物理约束与同一结算口径的前提下，仅屏蔽不同时刻的官方预报信息，开展8组信息消融实验，结果如表\ref{tab:q3_ablation}所示。
\begin{table}[htbp]
\centering
\caption{问题三不同日内预报组合的结果比较}
\label{tab:q3_ablation}
\scriptsize
\begin{tabular}{lccc}
\toprule
预报组合$\mathcal{S}$ & 年度费用/元 & 应急购电/kWh & 节省费用/元 \\ 
\midrule
$\varnothing$ & 15230376.12 & 334601.06 & 0 \\
$\{6\}$ & 15038617.00 & 266410.12 & 191759.11 \\
$\{12\}$ & 14746909.96 & 224631.56 & 483466.16 \\
$\{18\}$ & 14805047.54 & 228882.13 & 425328.58 \\
$\{6,12\}$ & 14679898.82 & 217968.28 & 550477.30 \\
$\{6,18\}$ & 14864586.68 & 218423.74 & 365789.44 \\
$\{12,18\}$ & 14710994.09 & 212056.29 & 519382.03 \\
$\{6,12,18\}$ & 14664345.40 & 210761.48 & 566030.72 \\
\bottomrule
\end{tabular}
\end{table}
由表\ref{tab:q3_ablation}可得，不使用任何日内新预报的方案费用为15230376.12元，全部启用三个日内预报后降至14664345.40元，节省56.60万元（3.72\%）。6:00、12:00、18:00更新的平均边际收益分别约为6.15、28.41、16.27万元，表明预报信息价值存在明显时点差异，其中12:00时点临近负荷晚高峰，信息经济价值最高。只采用任一单一时点费用均低于不使用日内预报，说明在题目提供的预报范围内三个日内时点均有实际作用。

\subsubsection{盈亏平衡更新成本}
将消融实验的费用差折算为单次合同更新的盈亏平衡成本：
\begin{equation}
\kappa^{*}=\frac{C_{\varnothing}-C_{\{6,12,18\}}}{\Delta n}=\frac{566030.72}{1002}\approx564.90\ \text{元/次},
\label{eq:q3_kappa}
\end{equation}
式中$\Delta n=1002$为全年合同更新次数。其经济含义为：若一次预报获取、传输、计算、合同重签与执行的综合成本低于564.90元，则全年启用全部日内预报在当前样本内具有正净收益。若只能选择一个时点，则12:00优先级最高。

\subsubsection{其他预报时点的近似测试}
题目仅在0:00、6:00、12:00、18:00提供独立官方预报；对其他时点的检验仅作为额外重优化机会的敏感性分析，不解释为新增独立预报的信息价值。近似测试显示，3:00和9:00的合同重算会增加现金费用；15:00的收益很小，仅可作为谨慎候选。21:00虽有数值改善，但此后实际光伏功率恒为零，其收益来自多一次基于当前SOC的合同重优化而非光伏预报信息，故不纳入推荐策略。

\subsubsection{代表日分析}
选取春、夏、秋、冬4个典型代表日，绘制日前计划购电$p_t$与最终合同$q_t$的轨迹，并标注6:00、12:00、18:00更新事件。不同季节光伏工况下，模型均会依据新预报对尚未执行的合同作向上或向下修正；指定时段的完整数值结果见附录。
\begin{figure}[htbp]
\centering
\includegraphics[width=0.85\textwidth]{figures/q3-contract.png}
\caption{问题三事件驱动的合同调整过程}
\label{fig:q3_contract}
\end{figure}

\subsection{本章小结}
本章建立支持日内购电合同修正的微网滚动MPC模型，并以不使用日内新预报的方案为基线完成334天仿真。模型在6:00、12:00、18:00基于当前SOC重优化剩余合同，最终合同只相对0:00原计划结算一次；分层风险裕量贯穿日前规划、合同调整与滚动执行。主策略相对基线使年度总费用下降3.72\%、应急购电量下降37\%。8组消融实验表明全部启用三个官方预报节省566030.72元，12:00的边际价值最高，单次合同更新的盈亏平衡成本约564.90元；补充时点测试不改变仅采用官方四个发布时点的策略选择。
